\begin{abstract}
Nonnegative aggregations of quadratic inequalities can certify that the
convex hull of a feasible set is proper, describe the hull exactly, or
approximate it with finitely many inequalities. We
separate these uses and resolve three conjectures of Blekherman, Dey,
and Sun, two affirmatively and one negatively. First, a
nonempty strict quadratic system with hidden hyperplane convexity (HHC),
or with a weaker asymptotic hyperplane condition, has a proper convex hull if and only if it admits
a nonconstant globally convex aggregation. The proof includes
the previously open case in which convex aggregations exist but all are
constant. Second, the Gram
map $(X,t)\mapsto(XX\transpose,t^2)$ has HHC exactly when $X$ has at
least as many columns as rows. This yields systems of three
inequalities in $2r\geq4$ variables with HHC whose open hull is not a
countable conjunction of strict quadratic inequalities, even ones
unrelated to the original system, and whose closed hull is not a finite
conjunction of nonstrict quadratic inequalities in the original
variables. Both hulls have finite semidefinite lifts, and
$N$ good aggregations approximate the closed hull with best Hausdorff
error $\Theta(N^{-2})$, with constants independent of $r$.
Third, for three inequalities whose homogenized matrices have a positive
definite real linear combination, at most four good aggregations
describe every nonempty proper strict hull, and four is sharp in
dimension at least three; the proof transfers the four-bound of
Blekherman and Dunbar for regular nonstrict systems to arbitrary strict
systems. Lean proofs cover the certificate theorem,
specified consequences, and the infinite-aggregation example with its
exact hulls and approximation rate.
\end{abstract}

\section{Introduction}
\label{sec:introduction}

Quadratic inequalities are a basic source of nonconvexity in
optimization. A nonnegative weighted sum of the inequalities is valid and
can reveal convex structure that no single row has; such aggregations
underlie convexification results and semidefinite relaxation
certificates. We ask whether some aggregation certifies that the convex
hull is not the whole space, whether aggregations describe the hull
exactly, and how well finitely many approximate it. The answers differ
sharply.

\paragraph{Prior work.}
Dines proved that the joint image of two homogeneous quadratic forms is
convex \citep{Dines1941}. Convexity of such images is the geometric basis
of the S-lemma and its extensions; Polyak proved it for three forms with
a positive definite real linear combination in dimension at least three
\citep{Polyak1998}. For convex hulls of quadratic sets, Yildiran described the two-inequality case using at most
two aggregations \citep{Yildiran2009}, and Dey, Mu\~noz, and Serrano
obtained aggregation descriptions for three inequalities whose
homogenized matrices have a positive definite real linear combination
\citep{DMS2022}; that condition allows signed coefficients, unlike an
aggregation.

Blekherman, Dey, and Sun introduced hidden hyperplane convexity (HHC):
the image of the homogenized quadratic map is convex on every linear
hyperplane \citep{BDS2024,BDSpreprint}. Under HHC, nonemptiness,
properness, and their dimension assumption, the hull is the
intersection of all \emph{good aggregations}. A good aggregation
contains the hull in its strict sublevel set and has at most one
negative homogeneous eigenvalue; it need not be globally convex. They
posed three questions \citep[Conjectures~3.1--3.3]{BDSpreprint}:
whether HHC can still require infinitely many good aggregations; whether
six can be necessary for three inequalities whose homogenized matrices
satisfy positive definite linear combination (PDLC); and whether, under
HHC, a nonempty strict system has a proper hull only if some nonconstant
globally convex aggregation exists. They had proved this when no nonzero
convex aggregation exists, assuming only that the image of the quadratic
parts is convex, which HHC implies; they had proved the certificate
characterization without HHC for separable systems with $n\geq2$; and
they had proved an upper bound of six under PDLC. Locators refer to the May~29, 2023 arXiv version
\citep{BDSpreprint}. Blekherman and Dunbar later proved a bound of four for
regular nonstrict systems under PDLC \citep{BD2025,BDpreprint}.

\paragraph{Contributions.}
\begin{itemize}
\item \emph{Convex certificates} (Section~\ref{sec:certificate}).
  Theorem~\ref{thm:certificate} shows that a nonempty strict system
  with asymptotic hyperplane convexity (convex images on arbitrarily
  distant sweeping hyperplanes for each normal, implied by HHC) has a
  proper convex hull if and only if some nonnegative aggregation has a
  positive semidefinite quadratic part and a nonzero quadratic or linear
  part. Corollary~\ref{cor:hhc-certificate} therefore
  answers Conjecture~3.3 affirmatively, including the case, not covered
  by the earlier zero-cone criterion, in which nonzero convex
  aggregations exist but their quadratic and linear parts all cancel. The
  key step eliminates the constant term using a strictly feasible point
  and normalizes in the finitely generated cone of quadratic and linear
  coefficients, so degenerating separation multipliers, which may produce
  only a negative constant, still yield an actual nonzero convex
  aggregation.
  Consequences (Section~\ref{sec:consequences}) include an exact test by
  $2n+1$ semidefinite programs (Corollary~\ref{cor:finite-sdp}) and, under
  the same hypotheses, equivalence of a whole-space hull with a
  whole-space Shor projection (Corollary~\ref{cor:shor-hull}); the
  underlying closure relation between the Shor projection and convex
  aggregations is classical \citep{FujieKojima1997,KojimaTuncel2000}.
  Examples delimit the hyperplane condition, strict feasibility, and the
  closure needed in the Shor comparison
  (Sections~\ref{sec:consequences}--\ref{sec:hypotheses}).
\item \emph{A sharp Gram criterion and infinite aggregation}
  (Sections~\ref{sec:gram-hyperplanes}--\ref{sec:infinite-aggregation}).
  The Gram map $(X,t)\mapsto(XX\transpose,t^2)$ with
  $X\in\R^{k\times r}$ has HHC exactly when $r\geq k$
  (Theorem~\ref{thm:gram-hhc}). Its hyperplane images are expressed
  through the classical squared fidelity \citep{Uhlmann2000}.
  Consequently, for $r\geq2$, the three inequalities defining
  \[
   S_r=\{(u,v)\in\R^r\times\R^r:
           \|u\|<1,\ \|v\|<1,\ u\transpose v>1/2\}
  \]
  have HHC, yet every exact description of $\conv S_r$ by
  strict good aggregations contains a prescribed continuum of rays
  (Theorem~\ref{thm:indispensable-rays}). This answers Conjecture~3.1
  affirmatively. More strongly, no countable conjunction of arbitrary
  strict quadratic inequalities describes the open hull, and no finite
  conjunction of arbitrary nonstrict quadratic inequalities describes
  its closure, in the original variables
  (Theorem~\ref{thm:no-finite-quadratics}). Both hulls have explicit
  finite semidefinite lifts, and a two-point decomposition proves the
  strict hull formula already in four variables
  (Theorem~\ref{thm:ball-hulls}). Section~\ref{sec:infinite-aggregation}
  compares this with earlier infinite examples, which lack HHC
  \citep{DMS2022}, and with semidefinite convexification of quadratic
  matrix programs \citep{Beck2007,Beck2009,WangKilincKarzan2024}.
\item \emph{Finite approximation} (Section~\ref{sec:approximation}).
  For $S_r$, the best Hausdorff error of the closed hull using at most
  $N$ good aggregations is $\Theta(N^{-2})$, with explicit constants
  independent of $r$ (Theorem~\ref{thm:approximation-rate}). The lower
  bound allows all good multipliers, including interior ones; the upper
  bound has an exact rational version with coefficients of $O(\log N)$
  bits (Corollary~\ref{cor:rational-mesh}). The exponent two is familiar
  from approximation of univariate convex functions \citep{Rote1992};
  here it concerns a prescribed cone of quadratic cuts in arbitrary
  dimension. By convex duality, a single good aggregation chosen for one
  linear objective attains its exact optimal value
  (Proposition~\ref{prop:one-objective}), so the uniform lower bound is
  not a lower bound on optimization iterations.
\item \emph{Four aggregations under PDLC}
  (Section~\ref{sec:four-aggregation}). Under PDLC, four good aggregations
  describe every nonempty proper strict hull, in every dimension $n\geq1$ (Theorem~\ref{thm:strict-four}). This answers
  Conjecture~3.2 negatively. We transfer the four-bound of Blekherman and
  Dunbar for regular nonstrict systems \citep{BD2025,BDpreprint} by
  compact inward perturbations and a limit that controls the negative
  components; the original matrices may be dependent and the weak
  system need not be regular or bounded. The four-necessary example
  of Blekherman, Dey, and Sun gives sharpness for $n\geq3$. Dunbar's
  dissertation states a stronger regular nonstrict four-bound
  \citep{Dunbar2025Thesis}; our contribution is only the transfer to
  unrestricted strict systems.
\end{itemize}

To the best of our knowledge, these results give the stated resolutions
of the three conjectures, the sharp Gram criterion, and the results for
the explicit example; we make no general claim that aggregation,
infinite necessity, semidefinite lifting, or the inverse-square exponent
originates here.

\paragraph{Organization.}
Section~\ref{sec:setting} fixes the setting; the cited sections
contain the results. Section~\ref{sec:formal-overview} states the exact
scope of the Lean verification. Appendix~\ref{app:application} gives a
local-certificate construction, and Appendix~\ref{app:three-span} treats
many rows in a three-dimensional matrix span, including a counterexample
to a tempting extension of the two-aggregation case.
